\subsection{Dihedral Groups}

In these exercises, $D_{2n}$ has the usual presentation 
$D_{2n} = \gen{r, s \mid r^n = s^2 = 1, \ rs = sr\inv}$.

\begin{exercise} [1.2.1]
    Compute the order of each of the elements in the following groups:
    \begin{subexercises}
        \item $D_6$ \quad (b) $D_8$ \quad (c) $D_{10}$
    \end{subexercises}
\end{exercise}

\begin{sol}
    For any reflection $sr^k \in D_{2n}$, it is clear that $sr^k \neq 1$. Moreover,
    \[(sr^k)^2 = s(r^ks)r^k = s(sr^{-k})r^k = r^{-k}r^k = 1\]
    so that every reflection has order 2. We now compute the orders of the rotations in each group.
    \begin{subexercises}
        \item $D_6 = \{1, r, r^2, s, sr, sr^2\}$. Then $|1| = 1$ and $|r| = |r^2| = 3$.
        \item $D_8 = \{1, r, r^2, r^3, s, sr, sr^2, sr^3\}$. Then $|1| = 1, |r^2| = 2$ and $|r| = |r^3| = 4$.
        \item $D_{10} = \{1, r, r^2, r^3, r^4, s, sr, sr^2, sr^3, sr^4\}$. Then $|1| = 1$ and $|r| = |r^2| = |r^3| = |r^4| = 5$. \qh
    \end{subexercises}
\end{sol}

\begin{exercise}[1.2.2]
    Use the generators and relations above to show that if $x$ is any element of $D_{2n}$ which is not a power of $r$, then $rx = xr\inv$.
\end{exercise}

\begin{sol}
    If $x$ is not a power of $r$, then $x$ is a reflection, i.e., $x = sr^k$ for some $0 \leq k < n$. Then
    \[rx = r(sr^k) = (rs)r^k = (sr^{-1})r^k = s(r^{-1}r^k) = s r^{k - 1} = (sr^k) r^{-1} = x r^{-1} \qh\]
\end{sol}

\begin{exercise}[1.2.3]
    Use the generators and relations above to show that every element of $D_{2n}$ which is not a power of $r$ has order $2$. Deduce that $D_{2n}$ is generated by the two elements $s$ and $sr$, both of which have order $2$.

    \hint{cf. \exref{1.1.33}.}
\end{exercise}
 
\begin{sol}
    If an element is not a power of $r$, then it is a reflection. See \exref{1.2.1} as to why every reflection has order $2$. Moreover, for any $r^k, sr^k \in D_{2n}$ for $0 \leq k < n$, we may rewrite them as
    \[r^k = (s(sr))^k \quad \text{and} \quad sr^k = s(s(sr))^k,\]
    showing that every element of $D_{2n}$ can be expressed in terms of $s$ and $sr$.
\end{sol}

\begin{exercise}[1.2.4]
    If $n = 2k$ is even and $n \ge 4$, show that $z = r^k$ is an element of order $2$ which commutes with all elements of $D_{2n}$. Show also that $z$ is the only non-identity element of $D_{2n}$ which commutes with all elements of $D_{2n}$.

    \hint{cf. \exref{1.1.33}.}
\end{exercise}
   
\begin{sol}
    Since $k \neq 1$, then $r^k \neq 1$. Moreover, $(r^k)^2 = r^n = 1$, hence $|r^k| = 2$. Since $r^k$ commutes with other powers of $r$, consider $sr^j \in D_{2n}$ for $0 \leq j < n$. Then
    \[r^k(sr^j) = (r^ks)r^j = (sr^{-k})r^j = s(r^{-k}r^j) = s r^{j-k} = (sr^j) r^k,\]
    Now let $w \in D_{2n}$ such that $wx = xw$ for all $x \in D_{2n}$. If $w = r^m$ for $0 \leq m < n$, then commuting with $s$ gives $w = w\inv$, which implies $m = k$ by \exref{1.1.33}. If $w = sr^m$ instead, then commuting with $r$ implies $r = r\inv$, which cannot be true since $n \geq 4$.
\end{sol}

\begin{exercise}[1.2.5]
    If $n$ is odd and $n \geq 3$, show that the identity is the only element of $D_{2n}$ which commutes with all elements of $D_{2n}$.
    
    \hint{cf. \exref{1.1.33}.}
\end{exercise}

\begin{sol}
    Using the same argument as the previous exercise and the odd case in \exref{1.1.33}, then $w$ cannot be a rotation since $r^k \neq r^{-k}$. The reflection case is also the same as the previous exercise.
\end{sol}

\begin{exercise}[1.2.6]
    Let $x$ and $y$ be elements of order $2$ in any group $G$. Prove that if $t = xy$, then $tx = xt\inv$ (so that if $n = |xy| < \infty$ then $x, y$ satisfy the same relations in $G$ as $s, r$ do in $D_{2n}$).
\end{exercise}

\begin{sol}
    Since $|x| = |y| = 2$, then $x = x\inv$ and $y = y\inv$. Then
    \[tx = (xy)x = x(yx) = x(xy)\inv = x t\inv \qh\]
\end{sol}

\begin{exercise}[1.2.7]
    Show that $\gen{a, b \mid a^2 = b^2 = (ab)^n = 1}$ gives a presentation for $D_{2n}$ in terms of the two generators $a$ and $b$ of order $2$ computed in \exref{1.2.3} above.
    
    \hint{Show that the relations for $r$ and $s$ follow from the relations for $a$ and $b$ and conversely, the relations for $a$ and $b$ follow from those for $r$ and $s$.}
\end{exercise}

\begin{sol}
    If $a^2 = b^2 = (ab)^n = 1$, put $r = ab$ and $s = a$. Then
    \[r^n = (ab)^n = 1, \quad s^2 = a^2 = 1, \quad rs = (ab)a = a(ba) = a(b\inv a\inv) = a(ab)\inv = sr\inv.\]
    Conversely, put $a = s$ and $b = sr$. The calculations are similar.
\end{sol}

\begin{exercise}[1.2.8]
    Find the order of the \defref{cyclic subgroup} of $D_{2n}$ generated by $r$.
\end{exercise}

\begin{sol}
    Since $|r| = n$, the cyclic subgroup generated by $r$ has order $n$ by \exref{1.1.35}.
\end{sol}

\begin{exercise} [1.2.9] 
    Let $G$ be the group of rigid motions in $\r^3$ of a tetrahedron. Show that $|G| = 12$.
\end{exercise}

\begin{sol}
    Label the vertices of a tetrahedron 1 through 4. Vertex 1 has 4 choices, vertex 2 has 3 remaining choices, and vertices 3 and 4 are determined after. It follows that $4(3) = 12$ symmetries.
\end{sol}

\begin{exercise} [1.2.10]
    Let $G$ be the group of rigid motions in $\r^3$ of a cube. Show that $|G| = 24$.
\end{exercise}

\begin{sol}
    8 choices for vertex 1 and 3 adjacent vertices for vertex 2; $8(3) = 24$ symmetries.
\end{sol}

\begin{exercise} [1.2.11]
    Let $G$ be the group of rigid motions in $\r^3$ of an octahedron. Show that $|G| = 24$.
\end{exercise}

\begin{sol}
    6 choices for vertex 1 and 4 adjacent vertices for vertex 2; $6(4) = 24$ symmetries.
\end{sol}

\begin{exercise}[1.2.12]
    Let $G$ be the group of rigid motions in $\r^3$ of a dodecahedron. Show that $|G| = 60$.
\end{exercise}

\begin{sol}
    20 choices for vertex 1 and 3 adjacent vertices for vertex 2; $20(3) = 60$ symmetries.
\end{sol}

\begin{exercise}[1.2.13]
    Let $G$ be the group of rigid motions in $\r^3$ of an icosahedron. Show that $|G| = 60$.
\end{exercise}

\begin{sol}
    12 choices for vertex 1 and 5 adjacent vertices for vertex 2; $12(5) = 60$ symmetries.
\end{sol}

\begin{exercise}[1.2.14]
    Find a set of generators for $\z$.
\end{exercise}

\begin{sol}
    Any integer is of the form $n1$, so $\z = \gen 1$.
\end{sol}

\begin{exercise}[1.2.15]
    Find a set of generators and relations for $\intmod$.
\end{exercise}

\begin{sol}
    Each element of $\intmod$ is of the form $k\bar 1$ for $0 \leq k < n$. Then $\intmod = \gen{\bar 1 \mid n\bar 1 = \bar 0}$.
\end{sol}

\begin{exercise}[1.2.16]
    Show that the group $\langle x_1, y_1 \mid x_1^2 = y_1^2 = (x_1 y_1)^2 = 1 \rangle$ is the dihedral group $D_4$ (where $x_1$ may be replaced by the letter $r$ and $y_1$ by $s$).

    \hint{Show that the last relation is the same as $x_1 y_1 = y_1 x_1^{-1}$.}
\end{exercise}

\begin{sol}
    Note that $y_1^2 = 1$, then $y_1\inv = y_1$. Then
    \[(x_1y_1)^2 = 1 \iff x_1y_1x_1y_1 = 1 \iff x_1y_1x_1 = y_1\inv = y_1 \iff x_1y_1 = y_1x_1\inv\]
    Then the relations of $D_4$ are satisfied by $x_1$ and $y_1$ if we set $r = x_1$ and $s = y_1$.
\end{sol}

\begin{exercise}[1.2.17]
    Let $X_{2n}$ be the group whose presentation is displayed in (1.2).
    \[X_{2n} = \gen{x, y \mid x^n = y^2 = 1, xy = yx^2}\]
    \begin{subexercises}
        \item Show that if $n = 3k$, then $X_{2n}$ has order $6$, and it has the same generators and relations as $D_6$ when $x$ is replaced by $r$ and $y$ by $s$.
        \item Show that if $(3, n) = 1$, then $x$ satisfies the additional relation: $x = 1$. In this case deduce that $X_{2n}$ has order $2$.
        
        \hint{Use the facts that $x^n = 1$ and $x^3 = 1$.}
    \end{subexercises}
\end{exercise}

\begin{solenum}
    \item If $n = 3k$, then
    \[X_{2n} = X_{6k} = \gen{x, y \mid x^{3k} = y^2 = 1, xy = yx^2}\]
    As shown in the text, $x^3 = 1$. Assume the relations of $X_{6k}$ hold. Set $r = x$ and $s = y$, so that $r^3 = x^3 = 1$, $s^2 = y^2 = 1$, and
    \[rs = xy = yx^2 = sr^2 = sr\inv.\]
    Now suppose the relations of $D_6$ hold. Then
    \[x^n = r^{3k} = (r^3)^k = 1^k = 1, \quad y^2 = s^2 = 1\]
    and
    \[xy = rs = sr\inv = sr^2 = yx^2\]
    Thus, the two presentations are equivalent, and $|X_{6k}| = |D_6| = 6$.
    \item Since $(3, n) = 1$, there exist integers $a$ and $b$ such that $3a + nb = 1$. From the relations of $X_{2n}$, we have $x^3 = 1$ and $x^n = 1$. Then
    \[x = x^{3a + nb} = (x^3)^a (x^n)^b = 1^a 1^b = 1\]
    so that $X_{2n} = {1, y}$ where $y^2 = 1$. It follows that $|X_{2n}| = 2$.
\end{solenum}

\begin{exercise}[1.2.18]
    Let $Y$ be the group whose presentation is displayed in (1.3).
    \[Y = \gen{u, v \mid u^4 = v^3 = 1, uv = v^2u^2}.\]
    \begin{subexercises}
        \item Show that $v^2 = v^{-1}$.
        
        \hint{Use the relation $v^3 = 1$.}
        \item Show that $v$ commutes with $u^3$.
        
        \hint{Show that $v^2 u^3v = u^3$ by writing the left-hand side as $(v^2u^2)(uv)$ and using the relations to reduce this to the right-hand side. Then use part (a).}
        \item Show that $v$ commutes with $u$.
        
        \hint{Show that $u^9 = u$ and then use part (b).}
        \item Show that $uv = 1$.
        
        \hint{Use part (c) and the last relation.}
        \item Show that $u = 1$. Deduce that $v = 1$, and conclude that $Y = 1$.
        
        \hint{Use part (d) and the equation $u^4 v^3 = 1$.}
    \end{subexercises}
\end{exercise}

\begin{solenum}
    \item $v^3 = 1 \implies v^2 = v\inv$.
    \item From $uv = v^2u^2$, we have $vuv = u^2$ so that $vu = u^2v\inv = u^2v^2$. It follows that
    \[v^2u^3v = (v^2u^2)(uv) = (uv)(uv) = u(vu)v = u(u^2v^2)v = u^3.\]
    \item It follows that $u^9 = (u^4)^2u = u$. Then $uv = (u^3)^3v = v(u^3)^3 = vu$.
    \item $uv = v^2u^2 = u^2v^2$. Then $1 = uv$.
    \item $u^4v^3 = u(uv)^3 = u1^3 = 1$. Then $1v = 1$, and $u = v = 1$, so $Y = 1$.
\end{solenum}